\section{Problems 1--15}

In Exercises 1--15, plot the point given in polar coordinates and then give three different expressions for the point such that: 
\begin{enumerate}[a)]
    \item $r < 0$ and $0 \leq \theta \leq 2\pi$
    \item $r > 0$ and $\theta \leq 0$
    \item $r > 0$ and $\theta \geq 2\pi$
\end{enumerate}

\hrule \bigskip

\begin{theorem}[Equivalent Representations of Points in Polar Coordinates]
    Suppose $(r, \theta)$ and $(r', \theta')$ are polar coordinates where $r \neq 0$, $r' \neq 0$ and the angles are measured in radians. Then $(r, \theta)$ and $(r', \theta')$ determine the same point $P$ if and only if one of the following is true.
    
    \begin{itemize}
        \item $r' = r$ and $\theta' = \theta + 2\pi k$ for some integer $k$ 
        \item $r' = -r$ and $\theta' = \theta + (2k + 1)\pi$ for some integer $k$
    \end{itemize}
    
    All polar coordinates of the form $(0, \theta)$ represent the pole regardless of the value of $\theta$.
\end{theorem} \medskip \hrule \bigskip

    Before we get started let's breakdown exactly what we're doing here. With polar coordinates we're given a distance from the origin, "r", and a degree of rotation, "$\theta$". Much like cartesian coordinates, the order in which we follow these directions is arbitrary. You can either travel the radius first and then account for the angle, or rotate first and then travel the radius. \medskip
    
    This gets a little funny with a negative radius. You can either travel the radius, going left on the x axis, and then rotate. Or, you can rotate according to the angle and then just go backwards a distance of "r". Think of this like adjusting the angle of a car and then driving in reverse. It's roughly the same as turning your car around 180 degrees and driving forward. Both methods can get you to the same location. One's just a little (a lot) more dangerous and illegal than the other. 